\documentclass[pdftex,english,oribibl]{llncs}

%% Spracheinstellungen laden
\usepackage[english]{babel}

%% Schriftart in der Ausgabe/Eingabe
\usepackage[T1]{fontenc}
\usepackage{textcomp}
\usepackage[latin1]{inputenc}

%% Zitate
\usepackage[numbers]{natbib}
\bibliographystyle{abbrvnat}
%\bibliographystyle{dinat}
%\bibliographystyle{plainnat}
%\bibliographystyle{splncs}
%% Similar to option "sectionbib" but \refname instead of \bibname
\makeatletter
\renewcommand\bibsection{\section*{\refname\@mkboth{\MakeUppercase{\refname}}{\MakeUppercase{\refname}}}}
\makeatother

%% Index
%\usepackage{makeidx}
%\makeindex

%% PDF Einstellungen
% muss nach natbib geladen werden!
\usepackage{nameref}
\usepackage{varioref}
\usepackage[pdfusetitle,pdftex,colorlinks]{hyperref}
\hypersetup{pdfborder={0 0 0}}
\hypersetup{bookmarksdepth=3}
\hypersetup{bookmarksopen=true}
\hypersetup{bookmarksopenlevel=1}
\hypersetup{bookmarksnumbered=true}
\usepackage{color}
\hypersetup{colorlinks=false}

%\usepackage[section]{tocbibind}

\makeatletter
\gdef\@keywords{}
\def\keywords#1{\gdef\@keywords{#1}}
\gdef\@subtitle{}
\def\subtitle#1{\gdef\@subtitle{#1}}

%% modified from llncs
\renewenvironment{abstract}{%
  \list{}{\advance\topsep by0.35cm\relax\small%
          \leftmargin=1cm%
          \labelwidth=\z@%
          \listparindent=\z@%
          \itemindent\listparindent%
          \rightmargin\leftmargin}%
          \item[\hskip\labelsep\bfseries\abstractname]}{%
  \if!\@keywords!\else{\item[~]\item[\hskip\labelsep\bfseries\keywordname]\@keywords}\fi%
  \endlist}

\AtBeginDocument{%
  \if!\@subtitle!\else\hypersetup{pdfsubject={\@subtitle}}\fi
  \if!\@keywords!\else\hypersetup{pdfkeywords={\@keywords}}\fi
}
\makeatother

% llncs hyperref fix
\makeatletter
\providecommand*{\toclevel@author}{0}
\providecommand*{\toclevel@title}{0}
\makeatother

%% Grafiken
\usepackage[pdftex]{graphicx}
\DeclareGraphicsExtensions{.pdf,.jpg,.png}
\usepackage{subfigure}

%% Mathe
\usepackage{amsmath}
\usepackage{amssymb}

%% Listings
\usepackage{listings}
\lstset{escapechar=\%, frame=tb, basicstyle=\footnotesize}

%% Sonstiges
\newcommand{\TODO}[1]{\par\textcolor{red}{#1}\marginpar{\textcolor{red}{TODO}}}
\newcommand{\TODOX}[1]{\textcolor{red}{#1}\marginpar{\textcolor{red}{TODO}}}
\pagestyle{plain}

% Keine "Schusterjungen"
\clubpenalty = 10000
% Keine "Hurenkinder"
\widowpenalty = 10000 \displaywidowpenalty = 10000

%%%%%%%%%%%%%%%%%%%%%%%%%%%%%%%%%%%%%%%%%%%%%%%%%%%%%%%%%%%%%%%%%%%%%%%%%%%%%%%
%%% BEGIN DOCUMENT
%%%%%%%%%%%%%%%%%%%%%%%%%%%%%%%%%%%%%%%%%%%%%%%%%%%%%%%%%%%%%%%%%%%%%%%%%%%%%%%
\title{Simulation-based resilience prediction of microservice architectures}
% \subtitle{My (optional) Subtitle}
\author{Samuel Beck, Johannes G�nth�r, Christoph Zorn}
\institute{University of Stuttgart\\Reliable Software Systems (RSS)\\70569 Stuttgart, Germany}


\begin{document}

\maketitle

\begin{abstract}
	Current software simulators are tailored towards one specific purpose of conservative software simulation. Given the success of these tools it would be useful to run these tools on microservice architectures. This paper will focus on the development of a simulator that can be used for microservice architectures.
\end{abstract}

\section{Introduction}\label{sec:introduction}

	Was rein muss: Was wir wollen (unser ziel), ...
	
	%This is my introduction. %\citet{Shaw2003WritingGoodSoftwareEngineeringResearchPapersMinitutorial} wrote a paper with %hints on how to write good software engineering research papers. By the way, this was an %example for using the \textit{natbib} command \texttt{\textbackslash{}citet\{\}}.
	
	%Section~\ref{sec:anotherSection} presents everything one must know. The conclusions follow %in Section~\ref{sec:conclusions}.

%\section{Another Section}\label{sec:anotherSection}
	
	%\textit{AspectJ} can be used to weave cross-cutting concerns into Java programs %\citep{AspectJ2007}. By the way, this was an example for using the \textit{natbib} command %\texttt{\textbackslash{}citep\{\}}.
	
	%We will now demonstrate how to use subfigures (see Figure~\ref{fig:subfig}). %Figure~\ref{fig:circle} shows a circle. A star is displayed in Figure~\ref{fig:star}.
	
	%\begin{figure}
		% \centering
		%\subfigure[This is a circle.]{\label{fig:circle}
		 % \includegraphics[width=0.3\textwidth]{figures/template_circle.pdf}
		%}
		%\subfigure[This is a star.]{\label{fig:star}
		%  \includegraphics[width=0.3\textwidth]{figures/template_star.pdf}
		%}
		%\caption{A circle (a) and a star (b). Note that any caption ends with a full stop %character.}
		% \label{fig:subfig}
	%\end{figure}
  
\section{Research}\label{sec:research}

\subsection{Tools in Comparison}\label{subsec:toolsInComparison}

\subsubsection{Spigo}
	
	During our research on existing tools for microservice simulation we discovered a tool called spigo. It was written by Adrian Cockcroft a Amazon Web Services employee in the programming language go. Therefore the name spigo comes from Simulate Protocol Interactions in Go. \newline On the first view the tool looked very promising. Spigo contains a fairly simple JSON input and the structure of the parameters is intuitive. Each microservice architecture consists of multiple microservices. Each microservice has a name, a package inheritance a counter of the instances and dependencies to other microservices. The reason we like this tool is because it can simulate the occurrence of an error. Spigo uses the error monkeys from the simian army. But here lies already on of the biggest disadvantages. One can only simulate the failure of a single microservice during the execution. Another point against spigo is the weak documentation of the actual implementation as well as minor bugs that remain from 12 months ago. \newline In conclusion it is quite sure that spigo is not an option for us to use for an extension module. The restrictive and static code design makes it hard to follow the workflow and append another module.

\subsubsection{Palladio}

	Palladio is a software component model for business information systems to enable model-driven predictions on throughput, response time and resource utilization. \cite{Palladio2007} Our motivation to get information on the PCM was primarily to get knowledge from the inheritance of SimuLizar. A tool extended from the PCM and described later in this article. Component models provide many advanteages over object-oriented development approaches i.e. higher usability, quality and better test potential. \cite{Palladio2007}

\subsubsection{Simulizar}
	
	Simulizar is an extension of the Palladio Component Model. This tool was especially developed for systems that change at runtime i.e. cloudcomputing and virtualized infrastructure environments. Such dynamic system adapt to the environment in order to meet the quality-of-service requirements. In the most cases software systems rely on static and fixed resource management so only steady states can be predicted.\cite{Simulizar2013} Simulizar uses the so called MAPE-K feedback loop (Monitor, Analyze, Plan, Execute, Knowledge) which is used to react to changes that are done during runtime. \newline Palladio and especially Simulizer seem to be to complex and big to get an understanding of the tool in depth. The sheer amount of code and resources that is used for collecting data and making prediction would be just an overhead on the problems that we currently face.
	

\subsubsection{GreenCloud} 

	With this tool we tried a different approach. Greencloud was designed to calculate the energy consumption of datacenters\citep[P.1]{GreenCloud}. Knowing that this was kind of a long shot we had the idea to take a simulator that simulates distributed objects and map microservice abilities and requirements to these objects. Since the simulator was written to overlook datacenters and their components, a mapping would mean that the entire microservice system would be mapped to a datacenter in the current GreenCloud simulator. Instances of microservices would compare to a server\citep[P.2]{GreenCloud} that gathers metrics. These metrics are currently power consumption, CPU- utilization and workload \citep[P.3]{GreenCloud} and should be changed or replaced to throughput,workload and whatever metrics we require from a microservice. GreenCloud also includes connections\citep[P.3]{GreenCloud} between servers. The connection could possibly used to be a mapping of microservice connections, that we can interfere to implement chaos monkeys. \newline
	Sounding good in theory but taking a closer look at the tool some problems lead to a problems. The focus of GreenCloud is obviously the simulation of power consumption. The previous mentioned connection between servers is not actually modeled but just taken into consideration regarding the power consumption of switches that connect servers together. Additionally there is no sign of the possibility to scale Instances (servers) which is a major part in a microservice system. Lastly workloads are only specified as the computing power they require (e.g. 3 Million Instructions per Second)\citep[P.3]{GreenCloud}. Communication or a differentiation between machines(DB, Workstation, ...) takes additional resources and is currently not taken into account. This would make mapping to a microservice architecture incredibly hard.  
	In hindsight it seems to be not very likely that using GreenCloud would be a good idea. The differences between provided capabilities by GreenCloud and required capabilities by the new tool is just to large. 


\section{Simulatorspecification}\label{sec:specifictation}

\section{Simulatordocumentation}\label{sec:documentation}


\section{Conclusions}\label{sec:conclusions}

  These are my conclusions.

\bibliography{template}

\end{document}
